%! Polynomial Zeros and Factoring — Unit 3, Lesson 3.4 — Student Packet Cover Sheet
% The ONE place \namedateperiod belongs (Namestrip). One learning target per Learning Objective
% the lesson draws on (1.11.A; 1.5.A.1–2), plus the habit the crux is built on.
\documentclass[10pt]{article}
\usepackage{precalculus-article}
\usepackage{precalculus-boxes}
\usepackage{ltablex}
\keepXColumns

\begin{document}

% Full-bleed plum banner
\begin{tikzpicture}[remember picture, overlay]
  \fill[plum] (current page.north west)
    rectangle ([yshift=-1.05in]current page.north east);
\end{tikzpicture}
\vspace{-0.8in}

\begin{center}
  {\color{white}\textbf{\LARGE Precalculus}}\\[4pt]
  {\color{lilacmid}\large\textbf{Unit 3: Polynomial Functions}}\\[6pt]
  {\color{white}\textbf{\large Lesson 3.4 \quad Polynomial Zeros and Factoring}}
\end{center}

\vspace{0.3in}
\namedateperiod
\vspace{0.18in}

\begin{learningtargetbox}
\small
\begin{enumerate}[leftmargin=1.5em, itemsep=5pt]
  \item I can \textbf{factor a polynomial completely} by a strategy: the \textbf{GCF} first,
        then a \textbf{difference of squares}, a \textbf{sum or difference of cubes}, a
        trinomial or \textbf{quadratic form}, or \textbf{grouping} --- until no factor can be
        factored further over the real numbers.
  \item I can \textbf{solve a polynomial equation} by moving every term to one side,
        factoring, and using the \textbf{Zero Product Property} --- finding every solution,
        real and non-real.
  \item I can use a factored form to name a polynomial's \textbf{zeros}, their
        \textbf{multiplicities}, and its \textbf{$x$-intercepts}, and check that a degree-$n$
        equation has $n$ solutions, counting multiplicity.
  \item I can explain why dividing both sides of an equation by $x$ loses the solution
        $x = 0$, and solve the equation without losing it.
\end{enumerate}
\end{learningtargetbox}

\vspace{0.14in}

\begin{tocbox}
\small
\renewcommand{\arraystretch}{1.5}
\begin{tabularx}{\linewidth}{@{} c l X @{}}
\rowcolor{lilac}
\textbf{\#} & \textbf{Component} & \textbf{Description} \\
\midrule
1 & Warm-Up      & Spiral review of prerequisite skills \\
2 & Guided Notes & Factoring strategies, the Zero Product Property, real and non-real solutions \\
\rowcolor{goldbg}
3 & AP Practice  & \textbf{Extra credit, scored 1 to 5.} AP-style multiple choice and free response \\
4 & Homework     & \textbf{Graded.} DeltaMath assignment \textit{Lesson 3.4: Polynomial Zeros and Factoring} \quad Due: \blank{2.2cm} \\
\end{tabularx}
\end{tocbox}

\vspace{0.14in}

\begin{remindbox}
\small
A product is zero only when one of its factors is zero --- so a factored polynomial hands you
its zeros. \textbf{Factor first}, and \textbf{factor completely}: every factor you leave
unfinished hides solutions. A factor of $x$ is a solution too: $x = 0$. Some solutions are
\textbf{non-real} and come in conjugate pairs; they are zeros, but not $x$-intercepts. Counting
multiplicity, a degree-$n$ equation always has $n$ solutions --- use that count to check you
found them all.
\end{remindbox}

\end{document}
